% Generated draft: requires mathematical and source review.
\documentclass[11pt]{article}
\usepackage[margin=25mm]{geometry}
\usepackage{fontspec}
\setmainfont{Latin Modern Roman}
\usepackage{amsmath,amssymb,mathtools,unicode-math}
\setmathfont{Latin Modern Math}
\usepackage{enumitem,needspace,xurl,hyperref}
\setlength{\parindent}{0pt}
\setlength{\parskip}{5pt}
\setlength{\emergencystretch}{3em}
\newcommand{\N}{\mathbb{N}}
\newcommand{\Z}{\mathbb{Z}}
\newcommand{\Q}{\mathbb{Q}}
\newcommand{\R}{\mathbb{R}}
\newcommand{\C}{\mathbb{C}}
\newcommand{\F}{\mathbb{F}}
\newcommand{\A}{\mathbb{A}}
\newcommand{\PP}{\mathbb{P}}
\newcommand{\E}{\mathbb{E}}
\newcommand{\eps}{\varepsilon}
\newcommand{\dd}{\,\mathrm{d}}
\begin{document}
\section*{UnsolvedMath ID: 30000978}
\textbf{Source problem code:} catalogue entry
\subsection{Definitions and mathematical statement}
Let $n\geq1$, let $S$ be a K3 surface, and let $(F,g)$ be a polarized holomorphic symplectic manifold deformation-equivalent to the Hilbert scheme $S^{[n]}$ of length-$n$ subschemes of $S$. Thus $F$ has complex dimension $2n$ and $g$ is the given ample polarization. Let $h$ be a divisor class. Write $g\cdot R$ and $h\cdot R$ for divisor--curve pairings, and $(R,R)$ for the intended Beauville--Bogomolov square on curve classes, obtained from the form on cohomology by the prescribed dual convention.

The first conjectural criterion is
\[
h\text{ is ample}
\quad\Longleftrightarrow\quad
h\cdot R>0\quad\text{for every intended curve class }R
\text{ with }g\cdot R>0\text{ and }(R,R)\geq-\frac{n+3}{2}.
\]
The source does not specify whether the test classes are actual effective curves, integral algebraic homology classes, primitive classes, or some other prescribed set. The criterion is retained with this qualification; no extra positivity-cone hypothesis on $h$ is silently added.

The second conjectural criterion states that, for each integer $r$ with $1\leq r\leq n$, extremal rays with a representing class $R$ satisfying
\[
-\frac{r+3}{2}\leq(R,R)
\]
should correspond to contractions whose exceptional locus is birationally fibered by projective spaces $\mathbb P^j$, with codimension $j\leq r$, and for which $R$ is represented by a line in a fiber. In the usual interpretation, codimension is measured in $F$. The source does not specify the cone defining extremality, existence on $F$ versus a birational model, the exact birational-fibration meaning, or treatment of multiple exceptional components.

A ray alone has no fixed square: replacing $R$ by $aR$ multiplies its square by $a^2$. Thus the representing-class normalization and the asserted identification with a line class are essential missing conventions. No additional upper bound on $(R,R)$ is inserted, and both numerical constants and the full range $1\leq r\leq n$ are preserved.

The corrected-verified note says both conjectural criteria are fully present and damaged intersection and projective-space notation was restored. In contrast, the literature assessment checked 2026-08-21 says that extraction ended immediately after ``Conjecture 3'' and omitted the condition, using that absence to label the entry open. That explanation is stale relative to the supplied restored statements; it cannot serve as a verified current mathematical status assessment. The registry's open label is retained as metadata.
\subsection{Short English statement}
Retain the proposed ampleness test using curve-square bound $-(n+3)/2$, and the contraction prediction using bounds $-(r+3)/2$. The tested curve classes, normalization of a ray generator and birational contraction conventions require review; the source's missing-text assessment conflicts with its restored statements.
\subsection{Sources}
\begin{itemize}
\item Local UnsolvedMath registry, permanent ID 30000978, source code catalogue entry; supplied complete record, including both restored criteria and the corrected-verified note, alongside its conflicting dated missing-text assessment.
\item Classical Algebraic Geometry (2009), Oberwolfach Reports, DOI 10.4171/owr/2008/26, pp. 1425--1474: \url{https://doi.org/10.4171/owr/2008/26}. Proposed-year metadata gives 2008.
\end{itemize}
\end{document}
